\item En un experimento de dispersión de Rutherford se disparan partículas alfa con una energía cinética de $7.70\text{ MeV}$ hacia un núcleo de oro que permanece en reposo durante la colisión. Las partículas alfa se aproximan al núcleo de oro hasta una distancia de $29.5\text{ fm}$ antes de retroceder.
\begin{enumerate}
    \item Calcule la longitud de onda de De Broglie para la partícula alfa de $7.70\text{ MeV}$ y compárela con la distancia de máxima aproximación de $29.5\text{ fm}$.
    \item A partir de esta comparación, explique por qué es físicamente más apropiado tratar la partícula alfa en este experimento como una partícula clásica en lugar de una onda cuántica.
\end{enumerate}